\chapter{Integrabilidad y teoremas de convergencia}\label{cap:v2-05}
Para integrar funciones con signo necesitamos evitar la cancelación de dos infinitos. Para intercambiar límites e integrales necesitamos controlar de qué manera la masa puede concentrarse o escapar. Este capítulo desarrolla las lecturas 9, 10 y 12 de Valencia Bucio \cite{valencia2}; el diagnóstico de pruebas falsas de convergencia monótona sigue la propuesta didáctica de Maximenko \cite{maximenkoFalsas}.

\section{Funciones reales integrables}\label{lectura:v2-09}
\begin{definicion}\index{parte positiva}\index{parte negativa}\index{funcion@función!integrable}
Para una función real medible $f$, definimos $f^+=\max(f,0)$ y $f^-=\max(-f,0)$. Entonces $f=f^+-f^-$ y $|f|=f^++f^-$. Decimos que $f$ es integrable si $\int|f|<\infty$, y ponemos
\[\int f=\int f^+-\int f^-.\]
Esta resta está definida porque ambos términos son finitos. Si solamente uno de ellos es finito, puede definirse una integral extendida, pero $f$ no se llama integrable.
\end{definicion}
\begin{lema}\label{lema:v2-descomposicion-integral}
Si $u,v,r,s$ son no negativas integrables y $u-v=r-s$, entonces $\int u-\int v=\int r-\int s$.
\end{lema}
\begin{proof}
La igualdad puntual equivale a $u+s=r+v$. La aditividad para funciones no negativas da $\int u+\int s=\int r+\int v$. Como las cuatro integrales son finitas, podemos reordenar y restar.
\end{proof}
\begin{teorema}\label{teo:v2-integral-lineal}
Las funciones reales integrables forman un espacio vectorial. La integral es lineal, monótona y satisface $|\int f|\le\int|f|$. Funciones medibles iguales casi en todas partes son simultáneamente integrables y tienen la misma integral.
\end{teorema}
\begin{proof}
La desigualdad $|af+bg|\le|a||f|+|b||g|$ prueba integrabilidad. Para aditividad escriba $f+g=(f^++g^+)-(f^-+g^-)$ y use el lema anterior para comparar esta descomposición con las partes positiva y negativa de $f+g$. Luego aplique aditividad no negativa a cada sumando. Para $a>0$, $(af)^+=af^+$ y $(af)^-=af^-$; para $a<0$, $(af)^+=(-a)f^-$ y $(af)^-=(-a)f^+$; para $a=0$ la conclusión es inmediata. Esto prueba homogeneidad en todos los casos.

Si $f\le g$, $g-f\ge0$ y linealidad da $\int g-\int f=\int(g-f)\ge0$. Como $- |f|\le f\le|f|$, se obtiene la desigualdad del valor absoluto. Finalmente, la igualdad casi en todas partes de $f,g$ implica la de sus partes positivas, negativas y valores absolutos. La proposición \ref{prop:v2-integral-positiva} iguala las integrales correspondientes, lo que prueba simultáneamente la equivalencia de integrabilidad y la igualdad de integrales.
\end{proof}
Para $A$ medible escribimos $\int_Af\,d\mu=\int_Xf\ind A\,d\mu$. La integral en un conjunto se entiende, por tanto, como una integral en todo el espacio de una función anulada fuera de ese conjunto.

\begin{proposicion}[Markov y anulación de la integral]\label{prop:v2-markov}\index{Markov, desigualdad de}
Si $f\ge0$ es medible y $t>0$, entonces
\[t\mu\{f\ge t\}\le\int_{\{f\ge t\}}f\,d\mu\le\int_Xf\,d\mu.\]
En particular, $\int f=0$ si y solo si $f=0$ casi en todas partes.
\end{proposicion}
\begin{proof}
Integre la desigualdad $t\ind{\{f\ge t\}}\le f\ind{\{f\ge t\}}\le f$. Si $\int f=0$, cada conjunto $\{f\ge1/n\}$ es nulo. Su unión es $\{f>0\}$ y es nula por subaditividad. La recíproca es la invariancia de la integral al modificar una función medible en un conjunto nulo.
\end{proof}

\section{Convergencia monótona y el papel de los cuantificadores}\label{lectura:v2-10}
\begin{teorema}[Convergencia monótona]\label{teo:v2-mct}\index{convergencia!monotona@monótona}
Si $f_n\ge0$ son medibles y $f_n\uparrow f$ casi en todas partes, con $f$ medible, entonces $\int f_n\uparrow\int f$.
\end{teorema}
\begin{proof}
Primero suponga convergencia y monotonía en todo punto. La monotonía de la integral da $L=\lim_n\int f_n\le\int f$. Fije una simple no negativa $s\le f$ y un número $c\in(0,1)$. Los conjuntos
\[E_n=\{f_n\ge cs\}\]
son medibles, crecen y cubren $X$: donde $s=0$ todos los puntos pertenecen a todos los $E_n$; donde $s>0$, el límite $f\ge s>cs$ obliga a pertenecer desde algún índice. Como $cs\ind{E_n}\le f_n$, obtenemos $\int f_n\ge c\int s\ind{E_n}$. Escribir $s$ como suma finita de indicadores y usar continuidad por abajo de $\mu$ da $\int s\ind{E_n}\uparrow\int s$. Luego $L\ge c\int s$. Si $\int s=\infty$, esto ya fuerza $L=\infty$; si es finita, hacer $c\uparrow1$ da $L\ge\int s$. Finalmente tome el supremo sobre todas las simples $s\le f$.

En el caso casi en todas partes, reúna en un único conjunto medible nulo todos los puntos donde falla la monotonía de algún par consecutivo o donde falla la convergencia. Anule $f_n$ y $f$ en ese conjunto. Las funciones obtenidas satisfacen el caso puntual y conservan sus integrales.
\end{proof}
\begin{observacion}[Dos pasos falsos y su reparación]
De $f_n\uparrow f$ no se deduce que algún $f_n$ alcance a $f$: las constantes $f_n=1-1/n$ nunca alcanzan el límite $1$. Tampoco de $\forall x\ \exists n:\ f_n(x)\ge cs(x)$ se deduce $\exists n\ \forall x:\ f_n(x)\ge cs(x)$. Por ejemplo, en $\N$, $f_n=\ind{\{1,\ldots,n\}}$ crece a $s=1$, pero ningún índice sirve para todos los puntos. La prueba correcta no busca ese índice común: utiliza los conjuntos crecientes $E_n$ y la continuidad de la medida. Este diagnóstico se inspira expresamente en \cite{maximenkoFalsas}; las pruebas deliberadamente falsas de esa presentación no se utilizan como demostraciones.
\end{observacion}
\begin{observacion}[La diagonal del argumento original]
También puede demostrarse el teorema tomando, para cada $j$, simples $s_{j,k}\uparrow f_j$ y definiendo $u_n=\max_{j\le n}s_{j,n}$. Se tiene $u_n\le f_n$, y $u_{n+1}\ge u_n$ porque cada fila es creciente y se añade una fila. Para cada $j$ fijo y $n\ge j$, $u_n\ge s_{j,n}$, luego $\lim_nu_n\ge f_j$; tomar supremo en $j$ da $\lim_nu_n\ge f$. La desigualdad inversa resulta de $u_n\le f_n\le f$. Así el teorema \ref{teo:v2-integral-aprox} también proporciona la conclusión y quedan explicadas las dependencias de índices de la lectura 10.
\end{observacion}

\begin{teorema}[Fatou]\label{teo:v2-fatou}\index{Fatou, lema de}
Si $f_n\ge0$ son medibles, entonces $\int\liminf_nf_n\le\liminf_n\int f_n$.
\end{teorema}
\begin{proof}
Defina $g_m=\inf_{n\ge m}f_n$. Son medibles, no negativas y crecen a $\liminf_nf_n$. Para todo $n\ge m$ se tiene $g_m\le f_n$, luego $\int g_m\le\inf_{n\ge m}\int f_n$. Convergencia monótona en el lado izquierdo y la definición de límite inferior en el derecho dan el resultado.
\end{proof}
\begin{teorema}[Convergencia dominada]\label{teo:v2-dct}\index{convergencia!dominada}
Sean $f_n$ medibles reales con $f_n\to f$ casi en todas partes, $f$ medible. Si existe $g\ge0$ integrable tal que $|f_n|\le g$ casi en todas partes para cada $n$, entonces $f_n,f$ son integrables,
\[\int|f_n-f|\longrightarrow0,\qquad\int f_n\longrightarrow\int f.\]
\end{teorema}
\begin{proof}
Fuera de la unión numerable de los conjuntos excepcionales, $|f|\le g$ y $|f_n-f|\le2g$. Por monotonía todas esas funciones son integrables. Aplique Fatou a $2g-|f_n-f|$, anulando previamente todas las funciones en el conjunto excepcional y donde $g$ es infinita. Su límite es $2g$, por lo que
\[2\int g\le\liminf_n\left(2\int g-\int|f_n-f|\right)
=2\int g-\limsup_n\int|f_n-f|.\]
Como $\int g$ es finita, la resta es legítima y obtenemos $\limsup_n\int|f_n-f|\le0$. La no negatividad implica convergencia a cero. Finalmente $|\int f_n-\int f|\le\int|f_n-f|$ prueba la convergencia de las integrales. La función dominante es una sola $g$; el subíndice accidental $g_n$ en el enunciado original no debe interpretarse como una dominación variable sin control adicional.
\end{proof}

\section{Funciones complejas e integración}\label{lectura:v2-12}
\begin{proposicion}\index{funcion@función!compleja medible}
Para $f:X\to\C$, son equivalentes: sus preimágenes de abiertos son medibles; sus preimágenes de borelianos son medibles; sus partes real e imaginaria son medibles.
\end{proposicion}
\begin{proof}
La familia de subconjuntos de $\C$ cuya preimagen es medible es una $\sigma$-álgebra; si contiene los abiertos, contiene los borelianos. Los niveles de las partes real e imaginaria son preimágenes de semiplanos abiertos, por lo que la primera condición implica la tercera. Recíprocamente, la preimagen de un rectángulo abierto $(a,b)+i(c,d)$ es $\{a<\Re f<b\}\cap\{c<\Im f<d\}$. Los rectángulos de extremos racionales forman una base numerable de $\C$; todo abierto es unión numerable de ellos. Esto prueba la tercera implicación y cierra las equivalencias.
\end{proof}
Las combinaciones lineales complejas y los productos de funciones medibles son medibles: sus partes reales e imaginarias son sumas y productos de funciones reales medibles. También lo es $|f|=\sqrt{(\Re f)^2+(\Im f)^2}$, pues la raíz cuadrada es continua.

\begin{definicion}
Una función compleja medible es integrable si sus partes real e imaginaria lo son, y se define $\int f=\int\Re f+i\int\Im f$.
\end{definicion}
\begin{teorema}\label{teo:v2-complejas}
La condición anterior equivale a $\int|f|<\infty$. La integral compleja es lineal y satisface $|\int f|\le\int|f|$. El teorema de convergencia dominada vale también para funciones complejas, incluyendo convergencia en integral del error absoluto.
\end{teorema}
\begin{proof}
Use $|\Re f|,|\Im f|\le|f|\le|\Re f|+|\Im f|$ para la equivalencia de integrabilidad. La aditividad se verifica por componentes. Si $\alpha=a+ib$ y $f=u+iv$, entonces $\alpha f=(au-bv)+i(av+bu)$; la linealidad real aplicada a las dos componentes da $\int\alpha f=\alpha\int f$.

Si $z=\int f=0$, la desigualdad del valor absoluto es inmediata. Si $z\ne0$, elija $\omega=\overline z/|z|$, de modo que $|\omega|=1$ y $\omega z=|z|$. Por linealidad,
\[\left|\int f\right|=\Re\int\omega f=\int\Re(\omega f)\le\int|f|.\]
Para convergencia dominada, $\Re f_n\to\Re f$ e $\Im f_n\to\Im f$ casi en todas partes y ambas sucesiones están dominadas en valor absoluto por $g$. El teorema real da convergencia de sus integrales. Además $|f_n-f|\to0$ y está dominada por $2g$; aplicarlo a esta sucesión real no negativa prueba la convergencia del error absoluto.
\end{proof}
El teorema de Fubini complejo de la lectura 12 se demostrará junto con su versión real en el capítulo \ref{cap:v2-06}; no se necesita para los resultados de este capítulo.

\section{Ejercicios sobre funciones integrables}
\begin{ejercicio}[Lectura 9, ejercicio 1]\label{ej:v2-l09-1}
Pruebe: (i) linealidad de la integral sobre $A$; (ii) integrabilidad sobre $A\cup B$ y $A\cap B$, e inclusión--exclusión, si $f$ es integrable sobre $A$ y $B$; (iii) monotonía de la integral sobre $A$.
\end{ejercicio}

\begin{ejercicio}[Lectura 9, ejercicio 2]\label{ej:v2-l09-2}
Construya una función no medible cuyo valor absoluto sea integrable.
\end{ejercicio}

\begin{ejercicio}[Lectura 9, ejercicio 3]\label{ej:v2-l09-3}
Demuestre la desigualdad de Markov y deduzca: (i) una función integrable es finita casi en todas partes; (ii) $\int|f|=0$ implica $f=0$ casi en todas partes; (iii) si $f$ es real integrable y $\int_Af\ge0$ para todo $A$ medible, entonces $f\ge0$ casi en todas partes.
\end{ejercicio}

\begin{ejercicio}[Lectura 9, ejercicio 4]\label{ej:v2-l09-4}
Dé dos funciones integrables cuyo producto no sea integrable.
\end{ejercicio}

\begin{ejercicio}[Lectura 9, ejercicio 5]\label{ej:v2-l09-5}
Si $f$ es integrable, $g$ es medible y $f\le g$ casi en todas partes, pruebe que la integral extendida de $g$ existe y $\int f\le\int g$.
\end{ejercicio}

\begin{ejercicio}[Lectura 9, ejercicio 6]\label{ej:v2-l09-6}\index{Riemann y Lebesgue, comparacion@Riemann y Lebesgue, comparación}
Si $f:[a,b]\to\R$ es Riemann-integrable, pruebe que es Lebesgue-integrable y ambas integrales coinciden.
\end{ejercicio}


\section{Ejercicios sobre intercambio de límites e integrales}
\begin{ejercicio}[Lectura 10, ejercicio 1: Beppo Levi]\label{ej:v2-l10-1}\index{Beppo Levi, teorema de}
Si $f_n\ge0$ son medibles, pruebe $\int\sum_nf_n=\sum_n\int f_n$.
\end{ejercicio}

\begin{ejercicio}[Lectura 10, ejercicio 2]\label{ej:v2-l10-2}
Sea $A\subseteq\R$ y una familia $(f_t)_{t\in A}$ de funciones integrables tal que $f_t(x)\to f_a(x)$ cuando $t\to a$ en $A$, para casi todo $x$ y cada $a\in A$. Si $|f_t|\le g$ casi en todas partes para una sola $g\in L^1$, pruebe la continuidad de $t\mapsto\int f_t$.
\end{ejercicio}

\begin{ejercicio}[Lectura 10, ejercicio 3: transformada de Laplace]\label{ej:v2-l10-3}\index{Laplace, transformada de}
Para $f$ Lebesgue-medible en $[0,\infty)$, defina $Lf(s)=\int_0^\infty e^{-sx}f(x)\,dx$ cuando la función es integrable. Pruebe: (i) su dominio es estable hacia la derecha; (ii) $Lf$ es diferenciable en el interior del dominio; (iii) si está definida en todo $\R$, $Lf(s)\to0$ cuando $s\to\infty$; (iv) aditividad.
\end{ejercicio}

\begin{ejercicio}[Lectura 10, ejercicio 4]\label{ej:v2-l10-4}
Si $|f_n|\le g$ casi en todas partes con $g$ integrable, pruebe integrabilidad de $\limsup f_n$ y $\int\limsup f_n\ge\limsup\int f_n$.
\end{ejercicio}

\begin{ejercicio}[Lectura 10, ejercicio 5]\label{ej:v2-l10-5}
Para $f_n=-n^{-1}\ind{[0,n]}$ en $\R$, pruebe convergencia uniforme a cero y fallo de la desigualdad de Fatou sin no negatividad.
\end{ejercicio}

\begin{ejercicio}[Lectura 10, ejercicio 6]\label{ej:v2-l10-6}
Si $f_n\ge-h$ con $h\ge0$ integrable, pruebe $\int\liminf f_n\le\liminf\int f_n$. Explique por qué no se aplica al ejercicio anterior.
\end{ejercicio}

\begin{ejercicio}[Lectura 10, ejercicio 7]\label{ej:v2-l10-7}
Extienda convergencia monótona a una sucesión creciente cuyo primer término es integrable, aunque no sea no negativo.
\end{ejercicio}

\begin{ejercicio}[Lectura 10, ejercicio 8: núcleos]\label{ej:v2-l10-8}\index{nucleo de medidas@núcleo de medidas}
Sea $K:X\times\Bcal\to[0,\infty]$ tal que $K(x,\cdot)$ es una medida en $(Y,\Bcal)$ y $K(\cdot,B)$ es medible para cada $B\in\Bcal$. Si $\mu$ es una medida en $X$ y $f:Y\to[0,\infty]$ es medible, pruebe: (i) $\nu(B)=\int_XK(x,B)\,d\mu(x)$ es una medida; (ii) $x\mapsto\int_Yf(y)K(x,dy)$ es medible; (iii) su integral en $X$ es $\int_Yf\,d\nu$.
\end{ejercicio}

\begin{ejercicio}[Lectura 10, ejercicio 9]\label{ej:v2-l10-9}
En el ejercicio anterior suponga $\mu(X)<\infty$, $M=\sup_xK(x,Y)<\infty$ y $f$ medible acotada real. Pruebe medibilidad y acotación de $x\mapsto\int_Yf\,K(x,dy)$ y la misma identidad de integrales.
\end{ejercicio}


\section{Una aplicación: derivación bajo la integral}
\begin{teorema}\label{teo:v2-derivar-integral}\index{derivacion bajo el signo integral@derivación bajo el signo integral}
Sea $F(t)=\int_Xh(x,t)\,d\mu(x)$ para $t$ en un intervalo abierto $I$. Suponga que cada $h(\cdot,t)$ es integrable, que $\partial_th(x,t)$ existe para todo $(x,t)$ y que $|\partial_th(x,t)|\le g(x)$ para una $g\in L^1$, $g\ge0$. Entonces $F$ es diferenciable y $F'(t)=\int_X\partial_th(x,t)\,d\mu(x)$.
\end{teorema}
\begin{proof}
Fije $t\in I$ y una sucesión $t_n\to t$, $t_n\ne t$. Los cocientes $q_n(x)=(h(x,t_n)-h(x,t))/(t_n-t)$ son medibles por las hipótesis de integrabilidad. Convergen puntualmente a $\partial_th(x,t)$, que por ello es medible. El teorema del valor medio aplicado a $s\mapsto h(x,s)$ da $|q_n(x)|\le g(x)$. No hace falta que el punto intermedio elegido por ese teorema dependa mediblemente de $x$, porque solamente usamos la cota. Convergencia dominada prueba
\[\frac{F(t_n)-F(t)}{t_n-t}=\int_Xq_n\,d\mu\longrightarrow\int_X\partial_th(x,t)\,d\mu.\]
El límite es el mismo para toda sucesión $t_n\to t$; el criterio secuencial para límites reales prueba la diferenciabilidad y la fórmula.
\end{proof}
